\section{Our Approach}
\label{sec:impl}

We implement five configurations that differ in how they maintain sorted data,
search for splits, and schedule tree nodes.  They share the learning semantics
and produce the same tree.  \ImplGpu{} and \ImplHybrid{} are two modes of the
same CUDA grower: the former disables CPU subtree hand-offs, whereas the latter
uses both devices.

\subsection{Shared Semantics and Data Model}

The input features are stored as column-major single-precision values, and
class labels are mapped to dense integers.  All configurations use the same
CART and C4.5 rule layer and the same split-scoring functions.  The CPU and
CUDA code compile the final gain calculation from a common header.  Gains
within \(10^{-12}\) are treated as ties; the earlier cut and then the lower
feature index win.  Filters used by the optimized implementations are
designed to discard only cuts that cannot replace the incumbent; surviving
candidates are scored by the shared double-precision calculation.

Except for \ImplSerialBase{}, the implementations represent each feature by a
presorted column of 8-byte entries containing the value and a packed row and
class identifier.  A node occupies the same contiguous range in every feature
column.  Stably partitioning these ranges creates two children whose columns
remain sorted, so a feature need not be sorted again.  Nodes created
concurrently are first held in a chunked store and finally renumbered into
flat preorder arrays; construction order therefore does not affect the output.
Tests require identical trees across all five configurations.  After growth,
the same CPU routines perform CART cost-complexity pruning or C4.5
post-processing and error-based pruning.

\subsection{\ImplSerialBase}

\ImplSerialBase{} is a direct sequential implementation and serves as an
executable reference.  At each node, it copies the node's row identifiers for
each feature, sorts them by feature value with a comparison sort, evaluates
every admissible cut, chooses the split, and recursively builds two new row
lists.  It thus repeats an \(O(Fn\log n)\) operation at every node and accesses
values and labels indirectly through row identifiers, while allocating
temporary vectors throughout growth.  It nevertheless invokes
the same stopping, split-selection, and pruning rules as the optimized
configurations, providing their correctness reference.

\subsection{\ImplSerialOpt}

\ImplSerialOpt{} replaces repeated sorting by a single-threaded stable radix
presort of every feature.  Thereafter, each split is a linear sweep followed
by a stable partition.  The winning feature is already divided by position:
its first \(n_L\) entries form the left child.  The other columns are
partitioned according to the corresponding row identifiers, preserving their
order.  Features constant within a node are removed for its whole subtree.

The sweep maintains inexpensive incremental estimates of the gain.  Exact
gain is computed only for candidates that can still win under conservative
rounding and convexity bounds; binary-class specialization and bulk processing
of equal-value runs reduce the remaining work.  The shared exact calculation,
in cut order, still determines the selected threshold.

Column storage is allocated once in slots.  The serial builder keeps one spare
slot: large ranges move out of place through a one-bit row map, while small
ranges use reusable local storage.  Child class counts are retained from the
winning sweep, and no partition is performed when both children are already
leaves.  Large allocations request transparent huge pages, and an explicit
stack replaces recursive growth.  \ImplParallel{} uses this same builder with
a thread pool enabled.

\subsection{\ImplParallel}

\ImplParallel{} combines two granularities of work.  For a node with at least
32,768 rows, feature sweeps and column partitions run in parallel.  After a
node with at least 4,096 rows is split, one child becomes a pool task while the
current worker continues with the other.  This uses data parallelism near the
root, where few large nodes exist, and task parallelism deeper in the tree.
The calling thread participates in nested parallel loops, and the grower waits
for all subtree tasks only at the end.

The parallel builder reserves up to
\(\min(F-1,\textit{threads})\) spare column slots so that several columns can
move concurrently.  Presorting and final preorder assembly are parallel as
well.  The CPU path has no socket-aware scheduler, per-socket data structure,
or explicit NUMA memory policy.  It respects the process affinity mask, while
Linux first-touch placement determines where pages reside; huge pages reduce
translation overhead but do not control placement.

\subsection{\ImplGpu}

\ImplGpu{} grows all nonterminal nodes on the GPU, breadth-first by tree level.
Each node--feature range is divided into tiles of 2,048 entries, exposing
parallelism across nodes, features, and values.  A level detects constant
ranges, forms tile histograms and prefixes, searches for the best cut of each
feature, and returns these compact results to the host.  The host applies the
same CART or C4.5 feature-choice rule as the CPU.  The device then stably
partitions the surviving columns and, while doing so, forms the histograms
needed by the next level.  Thus only one split-decision round trip is required
per level.

The device uses the same packed entries in one slot per feature plus one spare,
and CUB performs a stable radix presort as uploaded columns become available.
On the evaluated A100 and H100 GPUs, cuts are scored directly in double
precision.  On devices with slow double precision, an optional
single-precision pass conservatively filters candidates before the shared
double-precision calculation.

Transfers use page-locked staging buffers rather than pinning the complete
dataset.  This is the only explicit NUMA placement in the implementation: on
a multi-socket host, the staging buffers are allocated by a thread bound to
the GPU's PCI NUMA node, and dedicated copy threads are pinned to CPUs on that
node.  General training memory and workers are not bound this way.
\ImplGpu{} sets the hand-off threshold to one row, so every nonempty subtree
remains on the device; CPU threads still stage transfers and run shared
post-processing.

\subsection{\ImplHybrid}

\ImplHybrid{} retains the level-synchronous GPU pipeline for large nodes but
sends smaller children to the \ImplParallel{} builder.  The automatic
threshold is \(N/8192\), clamped to 1,024--2,048 rows for CART, and \(N/4096\),
clamped to 1,024--4,096 rows for C4.5; it can also be set explicitly.  These
empirical defaults reflect that a GPU level has substantial fixed overhead,
whereas the CPU pool can grow many small subtrees efficiently.

For a hand-off, the GPU gathers only columns that can still split the child
into a bounded device ring.  Two copy streams transfer the data through pinned
buffers into a pageable host arena, and a CUDA host callback submits CPU tasks
when their data arrive.  The GPU continues with the next large-node level
while those tasks run.  The arena is pre-touched by CPU threads to avoid page
faults during transfer, but it is not bound to the GPU's NUMA node.  CPU and
GPU branches allocate nodes from the same concurrent store and are assembled
in preorder after both finish.  The scheme therefore assigns individual nodes
of one tree by size; it neither assigns whole forest trees to devices nor
reduces the CPU to a controller.

\begin{figure}[t]
\centering
\includegraphics[width=\linewidth]{figs/approach-hybrid.pdf}
\caption{Size-aware growth in \ImplHybrid{}.  The GPU processes the large
nodes of a level while small children are gathered, copied through pinned
staging buffers into a host arena, and submitted to the CPU pool.  Host
callbacks overlap CPU subtree growth with the next GPU level; both paths
contribute to one tree store.}
\label{fig:approach-hybrid}
\end{figure}
